\documentclass[11pt,a4paper]{article}

\usepackage[utf8]{inputenc}
\usepackage[T1]{fontenc}
\usepackage[french]{babel}
\usepackage{geometry}
\usepackage{graphicx}
\usepackage{float}
\usepackage{booktabs}
\usepackage{hyperref}
\usepackage{enumitem}
\usepackage{xcolor}

\geometry{margin=2.4cm}
\hypersetup{colorlinks=true, linkcolor=black, urlcolor=blue}
\setlist{itemsep=0.25em, topsep=0.4em}
\graphicspath{{captures/}}

\newcommand{\capture}[2]{%
  \begin{figure}[H]
    \centering
    \includegraphics[width=0.92\textwidth]{#1}
    \caption{#2}
  \end{figure}
}

\begin{document}

\begin{center}
{\large Université Mundiapolis}\\[0.4em]
{\large Cloud computing}\\[1.2em]
{\LARGE\bfseries Conteneuriser un modèle iris,\\ puis compter ses prédictions avec Redis}\\[1.2em]
Houda Erimi\\[0.3em]
15 octobre 2026
\end{center}

\vspace{1em}

\section*{Introduction}

Ce travail se fait en deux temps. D'abord je mets un petit modèle de classification dans Docker : un conteneur l'entraîne, un autre l'utilise. Ensuite, pour le TP~4, je transforme cette prédiction en API et j'ajoute Redis. Redis ne sert pas à deviner la fleur. Il compte seulement combien de fois l'API a été appelée, et il se souvient de ce nombre tant que son volume existe.

Le modèle travaille sur le jeu iris : 150 fleurs, 4 mesures en centimètres, 3 espèces. Quand je ne donne pas de mesures, l'exemple utilisé est une setosa, \texttt{5.1, 3.5, 1.4, 0.2}.

\section{Mettre le modèle dans Docker}

\subsection{Deux conteneurs, un volume}

L'entraînement est dans \texttt{train/}. Le script apprend une forêt aléatoire et écrit \texttt{iris\_model.joblib} dans \texttt{/model}. Ce dossier est un volume : le fichier reste après l'arrêt du conteneur.

L'inférence est dans \texttt{infer/}. Elle ne réapprend rien. Elle ouvre le fichier déjà écrit et affiche l'espèce. Les deux services du premier Compose partagent le même volume, et l'inférence attend que l'entraînement soit terminé.

\capture{01-point-de-depart.png}{Le dépôt au départ : entraînement, inférence, et le Compose qui les enchaîne.}

\capture{02-predict-iris.png}{Le script de prédiction. Il charge le modèle dans le volume et affiche l'espèce.}

\subsection{La version naïve, puis la version allégée}

La première image part de \texttt{python:3.11}, copie les fichiers, puis installe les bibliothèques. Elle fonctionne, mais elle est lourde : presque toute la place est prise par l'image Python complète.

Je l'ai ensuite allégée de quatre façons.

\begin{itemize}
  \item La base devient \texttt{python:3.11-slim}.
  \item Le build a deux étapes. La première installe les paquets, la seconde ne garde que le résultat.
  \item \texttt{requirements.txt} est copié et installé avant le script Python. Un changement de code ne relance plus \texttt{pip}.
  \item Un \texttt{.dockerignore} laisse de côté les caches et les fichiers inutiles.
\end{itemize}

L'image finale a quand même besoin de \texttt{libgomp1}, sinon scikit-learn ne démarre pas.

\capture{03-dockerfile-inference.png}{Dockerfile d'inférence optimisé, avant d'en faire une API.}

J'ai mesuré les deux versions sur la même machine.

\begin{center}
\begin{tabular}{@{}lrrr@{}}
\toprule
Version & Taille & Build à froid & Build à chaud \\
\midrule
Naïve, \texttt{python:3.11} & 1,98 Go & 53 s et 76 s & environ 2 s \\
Optimisée, \texttt{python:3.11-slim} & 559 Mo & 65 s et 62 s & environ 2 s \\
\bottomrule
\end{tabular}
\end{center}

Les deux durées à froid sont l'entraînement puis l'inférence. L'image optimisée est environ 3,5 fois plus petite. Le premier build ne devient pas vraiment plus court, parce qu'il faut encore télécharger NumPy, SciPy et scikit-learn. Le gain, c'est la taille. Une fois les couches en cache, les deux versions se reconstruisent en deux secondes.

\subsection{Ce que Trivy a trouvé}

J'ai scanné l'image optimisée avec Trivy. Il n'y avait aucune faille critique. En revanche, deux failles hautes avaient déjà un correctif : \texttt{wheel} 0.45.1 (CVE-2026-24049) et \texttt{jaraco.context} 5.3.0 (CVE-2026-23949). Elles ne venaient pas de mon \texttt{requirements.txt}. Elles étaient embarquées dans \texttt{setuptools}, y compris dans des copies internes du paquet.

J'ai donc mis à jour les deux paquets, et \texttt{setuptools} lui-même en 84.0.0, sinon les vieilles copies restaient visibles pour le scan. Après reconstruction, Trivy ne signale plus de faille haute ni critique sur les paquets Python.

\section{TP 4 : l'API et Redis}

\subsection{Ce que renvoie l'API}

C'est le point que j'ai mis du temps à séparer. Une réponse de \texttt{/predict} contient deux choses qui n'ont pas le même rôle.

\begin{center}
\begin{tabular}{@{}lp{9.2cm}@{}}
\toprule
Champ & Rôle \\
\midrule
\texttt{mesures} & Les quatre tailles envoyées, ou l'exemple setosa. \\
\texttt{espece} & La prédiction. Le modèle choisit l'espèce. Ici, \texttt{setosa}. \\
\texttt{predictions} & Le compteur. Ce n'est pas une espèce. C'est le nombre d'appels à \texttt{/predict}. \\
\bottomrule
\end{tabular}
\end{center}

\texttt{/stats} ne calcule rien. Il redemande à Redis la valeur du compteur. Une réponse \texttt{\{"predictions":0\}} ne veut donc pas dire que le modèle s'est trompé. Elle veut dire qu'aucune prédiction n'a encore été demandée depuis le dernier effacement de Redis.

\subsection{Le code et le Compose}

\texttt{infer/api.py} charge le même fichier modèle. À chaque \texttt{GET /predict}, il devine l'espèce puis incrémente la clé \texttt{predictions} dans Redis. Le service Redis s'appelle \texttt{base}, comme demandé dans l'énoncé.

\capture{04-api-redis.png}{Les deux routes : prédire et compter, puis seulement relire le compteur.}

Le fichier \texttt{docker-compose.yml} lance les deux services. \texttt{api} publie le port 8000 et monte le volume du modèle déjà entraîné. \texttt{base} est \texttt{redis:7-alpine}, et son dossier \texttt{/data} est le volume \texttt{redis-data}. C'est ce volume qui garde le compteur.

\capture{05-compose-api-base.png}{Compose du TP : \texttt{api} et \texttt{base}.}

\texttt{docker compose config} ne démarre rien. Il réécrit la configuration telle que Compose la comprend. On y voit le port 8000, l'image Redis, et le volume externe du modèle.

\capture{06-compose-config.png}{Vérification avec \texttt{docker compose config}.}

\subsection{Le premier build a échoué}

Au premier \texttt{docker compose up --build}, Redis s'est bien téléchargé. L'installation Python, non. Le pip 24.0 livré avec l'image a refusé un paquet téléchargé, en disant que son hash ne correspondait pas. Le build s'est arrêté là.

\capture{07-build-erreur-1.png}{Le build s'arrête sur \texttt{pip install}.}

\capture{07-build-erreur-2.png}{Les paquets sont en train d'être téléchargés.}

\capture{07-build-erreur-3.png}{Le message de hash. Le nom du paquet n'est même pas indiqué.}

\capture{07-build-erreur-4.png}{Échec final du build, code 1.}

J'ai modifié le Dockerfile pour mettre pip à jour avant d'installer \texttt{requirements.txt}. Avec pip 26.2.1, la même installation passe. Au deuxième lancement, les couches sont en cache, Redis annonce qu'il accepte les connexions, et Uvicorn écoute sur le port 8000.

\capture{08-compose-up-1.png}{Deuxième lancement : l'image est construite, Redis démarre.}

\capture{08-compose-up-2.png}{L'API est prête. Les 404 viennent d'une sonde sur des adresses qui n'existent pas.}

Les lignes \texttt{404} sur \texttt{/} et \texttt{/json/version} ne viennent pas du TP. Quelque chose interroge ces adresses tout seul. Les seules adresses utiles sont \texttt{/predict} et \texttt{/stats}.

\subsection{Plusieurs prédictions}

J'ai d'abord ouvert \texttt{http://localhost:8000/predict} dans Firefox. La page restait blanche, avec seulement la case \og Pretty-print\fg{}. Pourtant le conteneur répondait \texttt{200}. Le JSON était bien renvoyé, le navigateur ne le dessinait pas.

\capture{09-navigateur-vide.png}{Firefox affiche une page vide alors que l'API a répondu.}

Avec \texttt{curl.exe}, la réponse est lisible. Deux appels de suite donnent la même espèce, setosa, et le compteur passe de 3 à 4. \texttt{/stats} confirme \texttt{\{"predictions":4\}}.

\capture{10-predict-stats.png}{Deux prédictions, puis le total lu dans Redis.}

\subsection{Le compteur et les volumes}

J'ai arrêté la pile avec \texttt{docker compose down}, puis je l'ai relancée avec \texttt{docker compose up}. Je n'ai pas rappelé \texttt{/predict}. \texttt{/stats} valait encore 4.

\capture{11-compteur-survit.png}{Après un simple \texttt{down}, le compteur est toujours à 4.}

Dans les logs de ce redémarrage, Redis écrit \texttt{Loading RDB} puis \texttt{keys loaded: 1}. Il a retrouvé sa sauvegarde sur le volume. \texttt{down} supprime les conteneurs, pas le volume.

J'ai ensuite fait \texttt{docker compose down -v}. Cette fois la sortie indique \texttt{Volume exercice-docker\_redis-data Removed}. Au \texttt{up} suivant, Redis démarre sans recharger de clé. \texttt{/stats} vaut 0.

\capture{12-compteur-perdu.png}{Après \texttt{down -v}, le compteur repart de zéro.}

C'est le résultat attendu des étapes 6 et 7. Le \texttt{0} concerne uniquement la mémoire du compteur. Le fichier du modèle est dans un autre volume, celui créé par l'entraînement, et il n'a pas été supprimé. Un nouvel appel à \texttt{/predict} répondrait encore \texttt{setosa}, et le compteur passerait à 1.

\section*{Conclusion}

Le modèle iris tourne dans deux images Docker, d'abord dans une version simple de 1,98 Go, puis dans une version de 559 Mo. L'entraînement et l'inférence se partagent le fichier du modèle par un volume. Le scan n'a pas montré de faille critique ; les deux failles hautes corrigées étaient des paquets trop anciens dans l'image de base.

Le TP~4 ajoute une API et Redis. L'API prédit l'espèce. Redis compte les appels. Ce compteur survit à \texttt{docker compose down}, parce que le volume reste. Il disparaît avec \texttt{docker compose down -v}, parce que le volume est supprimé avec les conteneurs. C'est cette différence que je retiens pour la suite, quand le même couple \texttt{api} et \texttt{base} sera déployé sur Kubernetes.

\end{document}
